\subsection{Actual canonical ratio metrics on the whole original normal model}

\begin{definition}[Original product and analytic restriction at every model point]
\label{local:actual-normalized-full-point-restriction}
\lean{EnvelopeProperty.ProjSubvariety.normalizedSimultaneousProductPointAt, EnvelopeProperty.ProjSubvariety.normalizedSimultaneousBaseRestrictionAt}
\leanok
\uses{local:actual-product-original-analytic-restriction}
Fix actual original ideals $I,J$ on the original product and their finite
original Gauss logarithms. Over any field, the SAME actual normalized
simultaneous blowup is normal and proper by the existing construction.
Every genuine relative valued point $y$ of this model maps through the
literal normalization and blowup projections to its ENTIRE original
product point, and then to its SAME original analytic variety restriction.
No choice of a Gauss lift of $y$ is involved.
\end{definition}
\begin{proof}
Compose the two actual scheme-valued point maps over the literal unchanged
polynomial base. Apply the whole original variety restriction already
constructed on all product points. The whole defining-ideal support
argument there ensures original variety membership.
\end{proof}

\begin{lemma}[True original pulled ideal logarithms at every model point]
\label{local:actual-normalized-full-point-ideal-log}
\lean{EnvelopeProperty.ProjSubvariety.modelIdealLog_normalizedSimultaneous_comap}
\leanok
\uses{local:actual-normalized-full-point-restriction}
For every actual original ideal $K$ on the original product, its double
geometric pullback through the SAME original blowup and normalization
has at $y$ the TRUE intrinsic logarithm of $K$ at the ENTIRE projected
original product point. Infinite logarithms are included.
\end{lemma}
\begin{proof}
Apply the genuine proper-center ideal pullback formula first to the
normalization, then to the blowup. Actual properness of the model,
blowup and original product is already derived, so the true centers and
residue-field evaluations commute with these literal maps.
\end{proof}

\begin{definition}[Same actual canonical generic ratio section and weight]
\label{local:actual-normalized-full-point-ratio-section}
\lean{EnvelopeProperty.ProjSubvariety.normalizedSimultaneousRatioSectionAt, EnvelopeProperty.ProjSubvariety.normalizedSimultaneousRatioWeightAt}
\leanok
\uses{local:actual-normalized-full-point-ideal-log}
At any original model point $y$ for which the actual first product ideal
has finite logarithm, construct the actual canonical generic section of
the SAME original Cartier ratio line and define its weight by its TRUE
negative-log model norm. No rational trivialization or metric equality is
supplied. This is the actual ratio line used for the original Gauss metric.
\end{definition}
\begin{proof}
The original finite-logarithm data give integrality of the original
simultaneous blowup. Both ideal factors are genuinely Cartier there and
remain Cartier under the actual dominant normalization. The preceding
true ideal-log pullback formula and the first product ideal's nonzero
value imply that its actual normalized pullback norm is nonzero.
Apply the canonical actual Cartier ratio-section construction: tensor
the first ideal's generic inverse-unit section with the actual dual
transpose of the second ideal inclusion. Use the already derived
invertibility of the SAME original ratio line and its genuine model norm.
\end{proof}

\begin{theorem}[True actual canonical weight equals original ideal-log difference]
\label{local:actual-normalized-full-point-ratio-weight}
\lean{EnvelopeProperty.ProjSubvariety.normalizedSimultaneousRatioWeightAt_eq}
\leanok
\uses{local:actual-normalized-full-point-ratio-section}
At every such actual original normal-model point, the TRUE weight of
this SAME canonical ratio section is exactly the first original product
ideal logarithm minus the second at the ENTIRE projected point.
\end{theorem}
\begin{proof}
The actual canonical Cartier ratio metric theorem gives the difference
of the two genuine normalized pullback ideal logarithms. Apply the true
whole ideal-log pullback equality separately to the two original ideals.
All Cartier properties and first-factor nonvanishing used in this
calculation were derived in the preceding construction.
\end{proof}
